\silentchapter{2 Ikkje alle posisjonar er like sannsynlege}{ikkje alle posisjonar er like sannsynlege}
\prop{2}{}{Ikkje alle posisjonar er like sannsynlege.\footnote{Sannsynet er systemtilstandens vekt; ingen tenestekonfigurasjon flyt fritt.}}
\prop{2.1}{d}{Eit \textbf{seleksjonstrykk} er ein faktor som endrar sannsynsfordelinga over formrommet.\footnote{Wright (1932)}\footnote{Prosessorkostnad, nettverkslatens og minneavgrensing deformerer topologien.}}
\prop{2.11}{t}{Han er ingen regel som dikterer éi form, men ein gradient som gjer visse posisjonar meir sannsynlege enn andre.\footnote{Gradienten tvingar fram mønsteret; han dikterer korkje den eksakte kodelinja eller grensesnittpikselen.}}
\prop{2.111}{t}{Eit seleksjonstrykk åleine avgjer ikkje ein posisjon.\footnote{Eitt isolert optimaliseringskrav genererer ingen stabil arkitektur.}}
\prop{2.1111}{t}{Det avgrensar kvar posisjonen ikkje kan vere.\footnote{Feilhandtering definerer yttergrensene for den moglege utføringsstien.}}
\prop{2.11111}{t}{Form er det som står att når seleksjonstrykka har utelukka det dei kan.\footnote{Arkitekturen er restverdien etter at optimaliseringskrava har barbert vekk dei ikkje-levedyktige alternativa.}}
\prop{2.11112}{t}{Ho er ein rest, inga konstruksjon.\footnote{Systemutvikling er ei reduktiv feilretting, inga additiv skaping.}}
\prop{2.111121}{t}{Eit trykk som ikkje utelukkar nokon region, er inkje trykk; det er ein tom kategori.\footnote{Ein kravspesifikasjon utan avgrensingar kompilerer til kva som helst.}}
\prop{2.111122}{t}{Eit trykk som utelukkar alle regionar unnateke éin, er inkje trykk; det er determinisme.\footnote{Eit system med berre éin mogleg tilstand ber inga informasjon.}}
\prop{2.111123}{t}{Dei observerte trykka ligg mellom desse ytterpunkta.\footnote{Heuristisk søk opererer alltid i rommet mellom absolutt fridom og absolutt tvang.}}
\prop{2.12}{a}{I kvar funksjonell klasse verkar alltid meir enn eitt seleksjonstrykk samstundes.\footnote{Eit distribuert system er fanga mellom kryssande trykk; konsensuskravet dreg mot latens, tilgjengekravet mot partisjonering.}}
\prop{2.121}{t}{Einsidige formforklaringar er alltid ufullstendige.\footnote{Å tilskrive suksessen til eit grensesnitt utelukkande til friksjonsløysinga neglisjerer det bakenforliggjande minnehierarkiet.}}
\prop{2.1211}{t}{Verka berre eitt trykk, ville alle former i klassen vere identiske eller tilfeldig fordelte langs éin akse.\footnote{Eit nevralt nettverk trent på éin tapsfunksjon utan regularisering kollapsar i overtilpassing.}}
\prop{2.12111}{t}{Den observerte ikkje-uniformiteten, med variasjon langs fleire uavhengige aksar, krev fleire uavhengige trykk.\footnote{Høgdimensjonal spreiing i kodelager vitnar om eit mangfald av motstridande avhengnader.}}
\prop{2.1212}{t}{At eit trykk er latent, tyder ikkje at det er fråverande.\footnote{Søppelinnsamlinga krev syklusar sjølv når heapen er halvtom.}}
\prop{2.12121}{o}{Trykk som ikkje varierer i det observerte utvalet, lèt seg knapt telje.\footnote{Nettverksprotokollar framstår usynlege inntil pakketapet overstig terskelen.}}
\prop{2.121211}{t}{Metodisk krev dette variasjon i vilkåra, ikkje berre i formene.\footnote{Stresstesting av mikrotenester krev nettverkspartisjonering, ikkje berre ulik last.}}
\prop{2.13}{a}{For minst eitt par av seleksjonstrykk peikar gradientane i motstridande retningar i minst éin region av formrommet.\footnote{CAP-teoremet formaliserer korleis konsistens og tilgjenge peikar i diametralt motsette retningar under partisjonering.}}
\prop{2.131}{t}{Sidan fleire uavhengige trykk dreg i ulike retningar, kan ikkje alle stettast optimalt samstundes.\footnote{Ingen relasjonsdatabase kan gje både maksimal normalisering og minimal leselatens på same tid.}}
\prop{2.1311}{t}{Kvar realisert form er eit kompromiss.\footnote{Michl (1995)}\footnote{Mikrotenestearkitekturen tolererer bandbreiddetap for å vinne isolasjon.}}
\prop{2.13111}{t}{Kompromisset er ingen veikskap; det er den einaste moglege balansen under rådande vilkår.\footnote{Eventuell konsistens er ingen systemfeil, men ei naudsynt tilpassing til eit asynkront nettverk.}}
\prop{2.13112}{t}{Sidan fleire gyldige kompromiss finst, er skilnaden mellom to former ingen feil, men eit uttrykk for at kreftene lèt seg balansere på meir enn éin måte.\footnote{Monolittiske og mikrokjernebaserte operativsystem balanserer nøyaktig dei same maskinvarekreftene ulikt.}}
\prop{2.13113}{t}{Å kalle ei form «suboptimal» føreset at ein veit kva ho er suboptimal i høve til.\footnote{Ein $O(n^2)$-algoritme er optimal dersom minneavgrensinga trumfar prosessortida.}}
\prop{2.131131}{t}{Utan eksplisitt spesifisering av trykksettet er «suboptimal» ikkje falsifiserbart.\footnote{Utan å vekte yting matematisk mot lesbarheit, er vurderinga av eit kodeparadigme uformell.}}
\prop{2.14}{t}{Funksjonen definerer klassen og med det formrommet.\footnote{API-kontrakten definerer tenesteklassen og stakar ut integrasjonsrommet.}}
\prop{2.141}{t}{Innanfor klassen er funksjonen konstant.\footnote{Forretningslogikken ber inga informasjon om kvifor to klientar med ulikt grensesnitt visualiserer tilstanden ulikt.}}
\prop{2.1411}{t}{Det konstante har null varians; det ber null informasjon om kvifor éi bestemt form vart realisert framfor ei anna.\footnote{Protokollstandarden forklarer ingenting om kvifor to rutarar implementerer bufferhandteringa ulikt.}}
\prop{2.14111}{t}{All observert formvariasjon under konstant funksjon er driven av dei resterande seleksjonstrykka.\footnote{Variasjonen i databasemotorar skuldast utelukkande I/O-kostnad, minnehierarki og trådhandtering, aldri SQL-standarden i seg sjølv.}}
\prop{2.141111}{i}{Stolen og sofaen har ulik funksjon og utgjer ulike klasser.\footnote{Køyretydsmiljøet og kompilatoren utgjer radikalt ulike klasser, skilt av tidsdomenet dei opererer i.}}
\prop{2.141112}{i}{To stolar med ulik stil har same funksjon. All skilnad mellom dei er ikkje-funksjonell.\footnote{To sorteringsalgoritmar oppnår eksakt same ordning. Avviket ligg utelukkande i rom- og tidskompleksitet.}}
\prop{2.141113}{o}{At materialet utgjer den mekaniske grensa, medan stilen manglar ergonomisk bodskap – og stilen likevel forklarer meir formvariasjon enn materialet – er ei direkte empirisk falsifisering av den funksjonalistiske doktrinen i denne klassen.\footnote{At programmeringsparadigmet forklarer meir av kodestrukturen enn prosessorarkitekturen, trass i at sistnemnde set dei fysiske grensene, motbeviser maskinens absolutte primat i programvareutvikling.}}
\prop{2.2}{d}{\textbf{Materialaffordansen} er operasjonane eit materiale tilbyr og yter motstand mot, under gjevne vilkår for kostnad og tilgjenge.\footnote{Gibson (1979)}\footnote{Affordansen i eit grafisk grensesnitt er ingen farge, men sekvensen av hendingar fargen logisk opnar for.}}
\prop{2.21}{t}{Signaturen er probabilistisk: Repertoaret bind sannsynsfordelinga, ikkje det einskilde objektet.\footnote{API-dokumentasjonen skildrar eit latent repertoar av moglege tilstandsovergangar, ikkje ein spesifikk sesjon.}}
\prop{2.211}{o}{Signaturen er latent.\footnote{Eit lukka system skjuler sitt fulle tilstandsrom bak den aktuelle tilstanden.}}
\prop{2.2111}{o}{Det teoretiske repertoaret slår ikkje tvingande ut i observerbar geometri.\footnote{Støtta for asynkron utføring i kjernen syner seg ikkje tvingande i sekvensiell brukarkode.}}
\prop{2.21111}{o}{Materialet er til stades, men algebraen manglar.\footnote{Kapasiteten for distribuert konsensus låg open i kryptografi og nettverksprotokollar i tiår før blokkjedene oppdaga syntaksen.}}
\prop{2.21112}{o}{Det som endrast over tid, er ikkje materialet, men kva operasjonar som er praktisk tilgjengelege.\footnote{Silisiumet er uendra, men samanstillinga av instruksjonssett har opna nye trajektoriar for maskinlæring.}}
\prop{2.211121}{i}{Latensen kan vere lang. Stål var tilgjengeleg i hundreår før den fyrste røyrstolen; materialet fanst, men operasjonen mangla.\footnote{Grafteori låg fast i eit århundre før sosiale nettverk realiserte relasjonsalgebraen kommersielt.}}
\prop{2.211122}{t}{Formhistoria er difor systematisk langsamare enn materialhistoria.\footnote{Systemarkitekturen vil alltid henge etter maskinvarens affordansar, av di operasjonane tek tid å formalisere til stabil grammatikk.}}
\prop{2.211123}{o}{Nye materiale arvar formspråket frå substratet dei erstattar, inntil eigne affordansar pressar dei inn i nye regionar av formrommet.\footnote{Dei fyrste webapplikasjonane imiterte trykte sider; fyrst seinare tvinga interaktiviteten fram asynkrone tilstandstre.}}
\prop{2.22}{t}{Ein samlevariabel, korrelert med alle samtidige seleksjonstrykk, bør predikere formvariasjon betre enn eitkvart einskildtrykk.\footnote{Designmønsteret fangar summen av fleire arkitektoniske trykk betre enn isolerte omsyn til minne eller prosessortid.}}
\prop{2.221}{o}{Stilperioden er ein slik samlevariabel.\footnote{Det programvaretekniske paradigmet fungerer presis slik; objektorientering kondenserer tiår med erfaring i modularisering og gjenbruk.}}
\prop{2.2211}{t}{Han er inga forklaring; han er ein proxy som absorberer alle samtidige trykk i éin etikett.\footnote{Termar som «Web 2.0» eller «skya» skildrar kvar vandringa stansa, men forklarer ingen av dei kausale drivkreftene.}}
\prop{2.22111}{t}{Stilkategoriane skildrar kvar vandringa stansa. Dei forklarer ikkje kvifor ho stansa nettopp der.\footnote{Å klassifisere eit system som «serverless» syner utfallet, ikkje optimaliseringsligninga som fann det.}}
\prop{2.22112}{t}{Å blande saman den deskriptive og den kausale krafta til stilomgrepet, er den vanlegaste feilen i formhistorieskrivinga.\footnote{Å tru at mikrotenester vann \textit{fordi} det var eit nytt paradigme, snur kausaliteten på hovudet; paradigmet er berre namnet på at dei vann.}}
\prop{2.221121}{t}{Stilperioden er symptomet; seleksjonstrykka er årsaka.\footnote{Arkitekturstilen er eit observerbart fenomen; distribusjonskostnad og vedlikehaldsbudsjett er dei verkelege kreftene.}}